\section{Containers, Images and the Registry}

\subsection{Why every job runs in a container}

A student's script is untrusted code on a shared PC. Running it directly
would let it read other students' files, fill the disk, or install software
that breaks the next job. A Docker container gives each attempt its own
filesystem view, process tree and resource limits, and the whole thing is
deleted afterwards.

\subsection{The hardened docker run}

\texttt{build\_docker\_run\_cmd()} in \texttt{agent/executor.py} builds this
command (abridged):

\begin{verbatim}
docker run --name lab-<attempt>
  --read-only                     root filesystem cannot be modified
  --network none                  no internet, no LAN, no exfiltration
  --pids-limit 512                fork bombs stop at 512 processes
  --cpus 1 --memory 512m          cgroup limits from the job's request
  --tmpfs /tmp:size=2g            scratch space in RAM, size-capped
  -v <input>:/input:ro            the student's code, read-only
  -v <output>:/output             the only writable place that is kept
  --user 10000:10000              not root (CPU jobs)
  --gpus all                      only when the job asked for VRAM
  <image> python /input/train.py
\end{verbatim}

The function also refuses to build any command that mounts
\texttt{docker.sock}: a container with the Docker socket can start a privileged
container and take over the PC. Each flag closes one door; together they make
a script able to compute and write results, and little else.

\paragraph{Cgroups.} \verb|--cpus| and \verb|--memory| are enforced by
Linux \emph{control groups}. A container over its memory limit is killed by the
kernel (exit code 137), which the agent reports as category \texttt{resource}.

\paragraph{GPUs in containers.} The NVIDIA Container Toolkit teaches Docker the
\verb|--gpus| flag: it mounts the GPU device and the host driver's libraries
into the container. Inside WSL2 the Windows driver is shared with Linux, which
is why lab PCs install the NVIDIA driver on Windows only.

\subsection{Tags and digests}

An image name like \texttt{pytorch/pytorch:2.5.1-cuda12.4-cudnn9-runtime} has
a \emph{tag}, which is a movable label: the publisher can point it at a
different image tomorrow. A \emph{digest} like
\texttt{registry:5000/lab-env/1bf18b82@sha256:2f7e3116...} is the SHA-256 of the
image manifest. Change one byte of the image and the digest changes, so pulling
by digest always gives exactly the same bytes.

The platform insists on digests for built environments: a build only counts
as ready when the builder reports a \texttt{repo/name@sha256:...} reference from
a registry. Two lab PCs running the same job therefore run the same software,
and a rerun next month runs what ran today. This is \emph{content addressing},
the same idea behind Git commit hashes and IPFS.

\subsection{Custom environments}

A student who needs \texttt{transformers==4.44.0} submits a
\texttt{requirements.txt}. Building it on every PC would be slow and could give
different results, so it is built \emph{once}:

\begin{enumerate}
\item The control plane normalises the requirements (sorted, lower-cased,
  duplicates removed; no \texttt{-r} or \verb|--index-url|) and hashes them
  together with the base image and GPU class:
\begin{lstlisting}[language=Python]
payload = {"requirements": reqs, "base": base_image, "compat": compat_class}
env_hash = sha256(json.dumps(payload, sort_keys=True).encode()).hexdigest()
\end{lstlisting}
\item The same hash means the same environment. A second student with the same
  requirements reuses the first student's image.
\item A \emph{builder} node (\verb|--role build|) writes a Dockerfile
  (\texttt{FROM base}, \texttt{RUN pip install -r requirements.txt}), builds it,
  runs an import test (and a CUDA test for GPU classes), records
  \texttt{pip freeze} as the lockfile, and pushes to the registry.
\item Jobs wait in \texttt{building} until the push succeeds, then every node
  pulls the image by digest.
\end{enumerate}

A broken \texttt{requirements.txt} fails on the builder in seconds, with the pip
error in the build log, instead of failing on a GPU node after a long wait.

\subsection{A bug: images tied to an old address}

The hash ignores the registry's address. When the laptop running the registry
changed Wi-Fi networks, its IP changed, but old builds still said ``ready'' with
digests pointing at the old IP. New jobs reused them and every pull failed. The
fix: when a ready build's digest names a different registry host than the
current \texttt{REGISTRY\_URL}, the build is reset to \texttt{pending} and
rebuilt, both at submit time and when a node claims. The lesson: when you cache
by content hash, ask what is \emph{not} in the hash.

\subsection{Docker-in-Docker for the simulated lab}

The simulated lab runs five ``PCs'' as containers on one laptop. Each runs its
\emph{own} Docker daemon inside (the \texttt{docker:dind} image, with
\verb|--privileged|). That matters: each simulated PC has its own image cache,
so a job on \texttt{sim-pc-3} really pulls from the registry like a separate
machine would, instead of sharing the laptop's cache.

\paragraph{Other use cases.} Container isolation is how GitHub Actions,
Google Cloud Run, AWS Lambda and Kubernetes run code; digests are how
Kubernetes and Helm pin exact images in production; build-once-pull-everywhere
is how every CI/CD pipeline ships software.
